\documentclass[a4paper,11pt]{article}

\usepackage[a4paper,margin=2cm]{geometry}
\usepackage[T1]{fontenc}
\usepackage[utf8]{inputenc}
\usepackage[ngerman,english]{babel}
\usepackage{lmodern}
\usepackage{microtype}
\usepackage{xcolor}
\usepackage{parskip}
\usepackage{enumitem}
\usepackage[most]{tcolorbox}
\usepackage{titlesec}
\usepackage[hidelinks]{hyperref}

\definecolor{HeaderBlueDark}{RGB}{70,110,170}
\definecolor{RowBlueLight}{RGB}{235,243,255}
\definecolor{RowBlueDark}{RGB}{220,233,250}
\definecolor{SoftGray}{RGB}{90,90,90}

\setlength{\parindent}{0pt}
\setlist[itemize]{leftmargin=1.4em,itemsep=0.35em,topsep=0.35em}
\linespread{1.06}

\titleformat{\section}[block]
{\Large\bfseries\color{HeaderBlueDark}}
{}{0pt}{}
\titlespacing{\section}{0pt}{1.3em}{0.8em}

\titleformat{\subsection}[block]
{\large\bfseries\color{HeaderBlueDark}}
{}{0pt}{}
\titlespacing{\subsection}{0pt}{1.0em}{0.6em}

\newtcolorbox{exbox}{enhanced,breakable,colback=RowBlueLight,colframe=RowBlueDark,boxrule=0.4pt,arc=1.5mm,left=1.2mm,right=1.2mm,top=1mm,bottom=1mm,title={Beispiele \textnormal{(Examples)}},fonttitle=\bfseries,colbacktitle=RowBlueDark,coltitle=black}

\NewDocumentEnvironment{grammarentry}{mm+b}{%
    \begin{tcolorbox}[enhanced,breakable,colback=white,colframe=HeaderBlueDark,boxrule=0.6pt,arc=1.5mm,left=1.5mm,right=1.5mm,top=1mm,bottom=1mm,colbacktitle=HeaderBlueDark,coltitle=white,fonttitle=\bfseries,title={#1\hfill\normalfont\footnotesize #2}]
    #3
    \end{tcolorbox}
}{}

\newcommand{\GermanText}[1]{\textbf{Deutsch: }#1\par}
\newcommand{\EnglishText}[1]{{\small\color{SoftGray}\textbf{English: }#1\par}}
\newcommand{\ExampleItem}[2]{\item \textbf{#1}\\{\small\color{SoftGray}#2}}

\title{Die Präposition \textit{ohne}}
\date{}

\begin{document}
    \maketitle
    \tableofcontents
    \newpage

\section{Grundbedeutung}

\begin{grammarentry}{ohne}{Präposition + Akkusativ}
\GermanText{\textit{ohne} bedeutet: nicht mit; etwas oder jemand fehlt.}
\EnglishText{\textit{ohne} means: without; something or someone is missing.}

\begin{itemize}
    \item \textbf{Kasus:} immer Akkusativ
    \item \textbf{Form:} ohne + Akkusativ
    \item \textbf{Englisch:} without
    \item \textbf{Gegenteil:} mit + Dativ
\end{itemize}
\end{grammarentry}

\section{ohne + Akkusativ}

\begin{grammarentry}{Kasus}{Akkusativ}
\GermanText{Nach \textit{ohne} steht das Nomen oder Pronomen im Akkusativ.}
\EnglishText{After \textit{ohne}, the noun or pronoun is in the accusative case.}

\begin{exbox}
\begin{itemize}
    \ExampleItem{Ich gehe ohne den Schlüssel aus dem Haus.}{I leave the house without the key.}
    \ExampleItem{Sie kommt ohne ihre Tasche.}{She comes without her bag.}
    \ExampleItem{Wir fahren ohne das Auto.}{We go without the car.}
    \ExampleItem{Er arbeitet ohne einen Plan.}{He works without a plan.}
\end{itemize}
\end{exbox}
\end{grammarentry}

\section{Artikel nach \textit{ohne}}

\begin{grammarentry}{Artikel}{Akkusativformen}
\GermanText{Weil \textit{ohne} Akkusativ verlangt, ändern sich besonders maskuline Artikel.}
\EnglishText{Because \textit{ohne} requires accusative, especially masculine articles change.}

\begin{center}
\begin{tabular}{lll}
\textbf{Genus} & \textbf{Nominativ} & \textbf{nach ohne} \\
Maskulin & der Schlüssel & ohne den Schlüssel \\
Feminin & die Tasche & ohne die Tasche \\
Neutral & das Auto & ohne das Auto \\
Plural & die Bücher & ohne die Bücher \\
\end{tabular}
\end{center}
\end{grammarentry}

\section{ohne + Pronomen}

\begin{grammarentry}{Pronomen}{Akkusativ}
\GermanText{Auch Pronomen stehen nach \textit{ohne} im Akkusativ.}
\EnglishText{Pronouns after \textit{ohne} are also in the accusative case.}

\begin{center}
\begin{tabular}{ll}
\textbf{Deutsch} & \textbf{English} \\
ich & ohne mich \\
du & ohne dich \\
er & ohne ihn \\
sie & ohne sie \\
es & ohne es \\
wir & ohne uns \\
ihr & ohne euch \\
sie/Sie & ohne sie/Sie \\
\end{tabular}
\end{center}

\begin{exbox}
\begin{itemize}
    \ExampleItem{Ohne mich kannst du nicht anfangen.}{You cannot start without me.}
    \ExampleItem{Ich mache das ohne dich.}{I do that without you.}
    \ExampleItem{Wir gehen nicht ohne ihn.}{We do not go without him.}
\end{itemize}
\end{exbox}
\end{grammarentry}

\section{ohne Artikel}

\begin{grammarentry}{Nullartikel}{without article}
\GermanText{\textit{ohne} kann auch direkt vor einem Nomen ohne Artikel stehen.}
\EnglishText{\textit{ohne} can also stand directly before a noun without an article.}

\begin{exbox}
\begin{itemize}
    \ExampleItem{Ich trinke Kaffee ohne Zucker.}{I drink coffee without sugar.}
    \ExampleItem{Er lebt ohne Angst.}{He lives without fear.}
    \ExampleItem{Sie arbeitet ohne Pause.}{She works without a break.}
    \ExampleItem{Ich kann ohne Hilfe nicht weitermachen.}{I cannot continue without help.}
\end{itemize}
\end{exbox}
\end{grammarentry}

\section{ohne zu + Infinitiv}

\begin{grammarentry}{Infinitivgruppe}{without doing something}
\GermanText{\textit{ohne zu + Infinitiv} bedeutet: jemand macht etwas nicht, während eine andere Handlung passiert.}
\EnglishText{\textit{ohne zu + infinitive} means: someone does not do something while another action happens.}

\begin{itemize}
    \item \textbf{Form:} ohne + zu + Infinitiv
    \item \textbf{Subjekt:} meistens dasselbe Subjekt wie im Hauptsatz
\end{itemize}

\begin{exbox}
\begin{itemize}
    \ExampleItem{Er geht weg, ohne etwas zu sagen.}{He leaves without saying anything.}
    \ExampleItem{Sie arbeitet, ohne eine Pause zu machen.}{She works without taking a break.}
    \ExampleItem{Ich habe geantwortet, ohne lange nachzudenken.}{I answered without thinking for a long time.}
\end{itemize}
\end{exbox}
\end{grammarentry}

\section{ohne dass + Nebensatz}

\begin{grammarentry}{Nebensatz}{without that}
\GermanText{\textit{ohne dass} leitet einen Nebensatz ein. Das konjugierte Verb steht am Ende.}
\EnglishText{\textit{ohne dass} introduces a subordinate clause. The conjugated verb is at the end.}

\begin{exbox}
\begin{itemize}
    \ExampleItem{Er ging weg, ohne dass ich es bemerkte.}{He left without me noticing it.}
    \ExampleItem{Sie hat geholfen, ohne dass jemand sie gefragt hat.}{She helped without anyone asking her.}
    \ExampleItem{Ich kann das nicht entscheiden, ohne dass du mir mehr Informationen gibst.}{I cannot decide that without you giving me more information.}
\end{itemize}
\end{exbox}
\end{grammarentry}

\section{Häufige Fehler}

\begin{grammarentry}{Fehler}{wrong / correct}
\GermanText{Nach \textit{ohne} kommt kein Dativ, sondern Akkusativ. Bei \textit{ohne zu} steht das Verb im Infinitiv.}
\EnglishText{After \textit{ohne}, use accusative, not dative. With \textit{ohne zu}, the verb is in the infinitive.}

\begin{exbox}
\begin{itemize}
    \ExampleItem{Falsch: Ich gehe ohne dem Schlüssel.}{Wrong: Dative after \textit{ohne}.}
    \ExampleItem{Richtig: Ich gehe ohne den Schlüssel.}{Correct: Accusative after \textit{ohne}.}
    \ExampleItem{Falsch: Er geht weg, ohne sagt etwas.}{Wrong: finite verb after \textit{ohne}.}
    \ExampleItem{Richtig: Er geht weg, ohne etwas zu sagen.}{Correct: \textit{ohne zu + Infinitiv}.}
\end{itemize}
\end{exbox}
\end{grammarentry}

\section{Kurze Zusammenfassung}

\begin{grammarentry}{Merken}{summary}
\begin{itemize}
    \item \textit{ohne} ist eine Präposition.
    \item \textit{ohne} verlangt immer Akkusativ.
    \item Grundbedeutung: \textit{without}.
    \item Gegenteil: \textit{mit + Dativ}.
    \item \textit{ohne zu + Infinitiv} benutzt man für Handlungen.
    \item \textit{ohne dass} leitet einen Nebensatz ein.
\end{itemize}
\end{grammarentry}

\end{document}
